% s3_method.tex
% Status: DRAFT — generated by TAR-Author; review before locking
% Lock condition: after ablation confirmed both components are needed

\section{Method: The Thermodynamic Continual Learning Governor}
\label{sec:method}

\subsection{Notation and Thermal Ratio}

Let $\theta \in \mathbb{R}^d$ denote the network parameters and let
$a^{(\ell)}_t \in \mathbb{R}^{B \times d_\ell}$ denote the activations of layer $\ell$
at step $t$ on a batch of size $B$.
Define the \emph{thermal level} as the mean activation standard deviation across layers:
\begin{equation}
    \sigma_t = \frac{1}{L}\sum_{\ell=1}^{L} \operatorname{std}\!\left(a^{(\ell)}_t\right).
\end{equation}
At the start of task $k$, after a warmup guard of $G = 60$ batches, we record the
\emph{task anchor}:
\begin{equation}
    \sigma^\star_k = \frac{1}{W}\sum_{t=G}^{G+W-1} \sigma_t,
    \quad W = 20 \text{ batches}.
    \label{eq:anchor}
\end{equation}
The anchor is frozen for the duration of task $k$ and reset at the start of task $k+1$.
The \emph{thermal ratio} is then:
\begin{equation}
    \rho_t = \frac{\sigma_t}{\sigma^\star_k}.
    \label{eq:rho}
\end{equation}
When $\rho_t > 1$ the network is more disordered than at the anchor moment; when
$\rho_t < 1$ it has consolidated below the anchor level.

\subsection{Warmup Guard}

The warmup guard delays anchor collection until the network has settled from the
activation distribution shift that occurs at each task boundary.
Without the guard, $\sigma^\star$ reflects random-initialisation noise rather than
the network's nominal thermal level, which makes the ordered regime ($\rho < 0.9$)
unreachable in practice --- a calibration bug identified during Phase~8 of this
investigation (see Appendix for the Phase~8 calibration report).

\subsection{Regime Classification and Learning Rate Modulation}

\begin{table}[h]
\centering
\caption{Thermal regime classification and governor actions.}
\label{tab:regimes}
\begin{tabular}{llll}
\toprule
$\rho_t$ & Regime & Interpretation & LR Action \\
\midrule
$> 1.1$ & \textbf{Disordered} & Activations more variable than anchor; exploring &
    $\eta \leftarrow \eta_0\!\left(1 + \alpha(\rho_t - 1)\right)$ \\
$[0.9, 1.1]$ & \textbf{Critical} & Near equilibrium & $\eta \leftarrow \eta_0$ \\
$< 0.9$ & \textbf{Ordered} & Activations below anchor; consolidating &
    $\eta \leftarrow \eta_0\,\alpha\,\rho_t$ \\
\bottomrule
\end{tabular}
\end{table}

The interpolation coefficient $\alpha = 0.5$ is pre-registered and controls how
aggressively the governor modulates the learning rate in response to thermal state.

\subsection{Dimensionality-Weighted L2 Anchor Penalty}

When the network enters the ordered regime on task $k$, we record the parameter
vector $\theta^\star_k = \theta$ as the task weight anchor.
For all subsequent tasks $k' > k$, we add the anchor regulariser:
\begin{equation}
    \mathcal{L}_{\text{anchor}} = \lambda_{\text{TCL}} \cdot D_{\mathrm{PR}} \cdot
        \|\theta - \theta^\star_k\|^2,
    \label{eq:penalty}
\end{equation}
where $\lambda_{\text{TCL}} = 0.01$ (pre-registered) and $D_{\mathrm{PR}}$ is the
\emph{dimensionality participation ratio} --- an estimate of the effective number of
active gradient directions.
In the experiments reported here, eigenvalue decomposition is disabled for
computational efficiency (\texttt{compute\_dpr=False}), fixing $D_{\mathrm{PR}} = 1.0$;
the thermal ratio computation is unaffected.

The anchor penalty differs from EWC in two ways.
First, it is triggered by the thermal signal rather than computed retrospectively from
a Fisher estimate after task completion.
Second, the $D_{\mathrm{PR}}$ factor scales the penalty to the effective dimensionality
of the current gradient subspace, providing stronger regularisation for compact
representations and weaker regularisation for diffuse ones.

\subsection{Full Governor Pipeline}

At each training step $t$ within task $k$:
\begin{enumerate}
    \item Compute $\sigma_t$ from the forward pass activations.
    \item If $t \geq G$ and the anchor has not yet been set, collect toward $\sigma^\star_k$
          (Eq.~\ref{eq:anchor}); freeze after $W$ batches.
    \item Compute $\rho_t = \sigma_t / \sigma^\star_k$ (Eq.~\ref{eq:rho}).
    \item Classify the regime (Table~\ref{tab:regimes}) and update $\eta$ accordingly.
    \item If ordered ($\rho_t < 0.9$) and $k > 1$, add $\mathcal{L}_{\text{anchor}}$
          (Eq.~\ref{eq:penalty}) to the training loss.
\end{enumerate}
The overhead relative to standard training is one scalar reduction per layer per step
(for $\sigma_t$) and one vector norm per step (for the penalty when active).
No Fisher matrix is computed; no per-parameter importance scores are stored beyond
the anchor parameter vector $\theta^\star_k$.

\subsection{Implementation}

TCL is implemented in \texttt{tar\_lab/thermoobserver.py}.
The governor wraps any PyTorch model and optimizer; experiment configuration is
managed via \texttt{ContinualLearningBenchmarkConfig} in \texttt{tar\_lab/schemas.py}.
All hyperparameters ($\lambda_{\text{TCL}} = 0.01$, $\alpha = 0.5$, $G = 60$,
$W = 20$) were pre-registered before Phase~10 data collection.
